\section{Linear Attention：把平方矩阵改写成线性递推}

本章解释 Linear Attention 的基本直觉。所谓 linear，不是模型变成线性函数，而是复杂度相对序列长度接近线性。典型做法是去掉或改写 Softmax，通过代数变换把注意力写成类似 RNN 的 recurrence，每一步维护一个状态，整体处理长度 \(L\) 的序列时成本接近 \(O(L)\)。

\begin{figure}[H]
\centering
\includegraphics[width=0.94\textwidth]{linear-attention-intuition.png}
\caption{Linear Attention 直觉：平方矩阵改写成线性 recurrence。自制概念图，依据 00:07:04--00:11:19 对谈内容整理。}
\end{figure}

\begin{knowledgebox}{读图：Linear Attention 的关键不是少写一层}
Softmax Attention 显式构造 \(L\times L\) 矩阵；Linear Attention 试图把历史信息压入递推状态，让每一步生成只更新状态。这样能降低复杂度和 KV cache 压力，但表达能力和训练稳定性要重新设计。
\end{knowledgebox}

\subsection{它在 Transformer 里的位置}

Transformer 主要由 Attention 和 FFN 反复堆叠构成。近几年，FFN 被 MoE 改造；Linear Attention 则是在 Attention 模块上动刀。它属于 pre-training 架构研究的一部分，和 optimizer、pre-training data、post-training 等不同。现在更常见的是 hybrid 架构：部分层保留 Full Attention，部分层换成 Linear Attention。

\begin{importantbox}{Hybrid 架构的动机}
纯 Linear Attention 可能表达力不足；纯 Full Attention 太贵。Hybrid 架构用少量全局层兜底表达能力，用大量线性层降低成本，是当前更实际的折中。
\end{importantbox}

\subsection{本章小结}

Linear Attention 的核心是复杂度和状态表示的重写。它能降低长上下文推理成本，但必须解决表达能力、训练稳定、硬件并行和大规模验证问题。

